\documentclass[10pt,a4paper]{amsart}
\usepackage[margin=19mm]{geometry}
\usepackage{amsmath,amssymb,mathtools}
\usepackage[scr=rsfs,cal=euler]{mathalfa}
\usepackage{xfrac}
\usepackage[unicode,hidelinks,bookmarks=false]{hyperref}
\newcommand{\ie}{i.e.\ }
\newcommand{\cf}{cf.\ }
\newcommand{\resp}{respectively\ }
\newcommand{\Subsection}{\subsection}
\providecommand{\nobreakdash}{\nobreak\hbox{-}}
\setlength{\emergencystretch}{2em}
\multlinegap0pt
\usepackage{bm}
\usepackage{bbm}
\usepackage{dsfont}
\usepackage{xspace}
\usepackage{cancel}
\usepackage{mathtools}
\usepackage{amsthm}
\makeatletter
\@ifundefined{theorem}{\newtheorem{theorem}{Theorem}}{}
\@ifundefined{lemma}{\newtheorem{lemma}[theorem]{Lemma}}{}
\makeatother
\makeatletter
\expandafter\ifx\csname customthm\endcsname\relax
\newenvironment{customthm}[1]
  {\renewcommand\theinnercustomthm{#1}\innercustomthm}
  {\endinnercustomthm}
\fi
\makeatother
\title[Infinitesimal Dilogarithm Satisfies Cluster Identities]{Infinitesimal Dilogarithm Satisfies Cluster Identities}
\author{Sinan Unver}
\date{Source: arXiv:2510.00016v1 (CC BY 4.0)}
\begin{document}
\maketitle

\noindent\emph{Source and licence.} Infinitesimal Dilogarithm Satisfies Cluster Identities. Version arXiv:2510.00016v1 (CC BY 4.0), distributed under the Creative Commons Attribution 4.0 International licence, as recorded in the arXiv licence metadata for this exact version and shown on the abstract page https://arxiv.org/abs/2510.00016v1. Full rights notice: licenses/arxiv-2510.00016-CC-BY-4.0.txt.\par\smallskip

\noindent\textit{Editorial scope.} Counted unit: the Theorem whose source label is \texttt{theorem-0-identity} in the article (source0.tex lines 311--317, source label \texttt{theorem-0-identity}), delivered together with the article's own complete original proof of it (source0.tex lines 320--361, one \texttt{proof} environment of 42 lines). No mathematics is written, repaired or re-derived: statement and proof are the article's own text line for line. Required context retained uncounted: dilog-formula (lines 232--234); periodic-sequence (lines 262--264); lemma (lines 290--296). Every definition, display and label of those lines is retained unchanged. Omitted from this package and not redistributed: 1--231, 235--261, 265--289, 297--310, 318--319, 362--578. Nothing inside a retained excerpt is omitted or altered: every retained body is delivered exactly as the article's lines are written, so no omission receipt of any kind accompanies this package. Citations: the retained excerpts carry no citation command at all, so no bibliography entry of the article is transcribed and no reference is redistributed. Reference adaptation: none was needed and none was made; every reference inside the delivered unit resolves inside it, so no unresolved reference or citation is produced. Printed slips kept literal: no printed word, symbol, hypothesis or number of the counted statement or of its proof was altered. Layout and class: the article's own class is \texttt{amsart} with packages including amscd, amsmath, amssymb, amsfonts, verbatim, xcolor, graphicx, xy, calrsfs; the wrapper is the lane's pinned \texttt{amsart} editorial wrapper at 10pt on A4, which restates the theorem-like environments the retained text needs and adds the article's own macro definitions that the retained text needs (none), with extra wrapper packages (bm, bbm, dsfont, xspace, cancel, mathtools). The delivered numbering is the unit's own. Assets: no retained span contains a figure environment, a graphics-inclusion command, a PDF-page inclusion, a table or a TikZ picture; no image file is copied into this package, the package declares no asset dependency and no image file is redistributed with it. Licence: version arXiv:2510.00016v1 of this article is distributed under the Creative Commons Attribution 4.0 International licence, as recorded in the arXiv licence metadata for this exact version and shown on the abstract page https://arxiv.org/abs/2510.00016v1. A full rights notice is delivered at licenses/arxiv-2510.00016-CC-BY-4.0.txt. Characters: every retained excerpt is pure ASCII, so the delivered engine, XeLaTeX with its default Latin Modern text font, reports no missing character.
\begin{eqnarray}\label{dilog-formula}
 g_{m,w} ([\tilde{\alpha}]) =\sum _{1\leq i \leq w-m} i \cdot (\ell _{w-i} \wedge \ell _{i})(\delta(\tilde{\alpha}))   
\end{eqnarray}

\begin{eqnarray}\label{periodic-sequence}
\Upsilon[0] \xrightarrow{r_0} \Upsilon[1] \xrightarrow{r_1} \cdots \xrightarrow{r_{P-1}} \Upsilon[P],   
\end{eqnarray}

\begin{lemma}\label{lemma}
Suppose that we are given a $\nu$-periodic sequence of mutations in a cluster pattern as in (\ref{periodic-sequence}). Let $k$ be a field with   ${\rm char}(k)\neq$2. There is a proper algebraic set $X \subseteq \mathbb{A}^{n}_{k}$ inside the $n$-dimensional affine space $\mathbb{A}^{n}_{k}$ over $k$ such that for ${\alpha}_{1}, \cdots, {\alpha} _{n} \in k_{\infty} $ with $({\alpha}_{1}(0), \cdots, {\alpha} _{n} (0)) \in (\mathbb{A}^n \setminus X)(k),$ we have 
$$
\sum_{0\leq j <P}\theta _{r_{j}}\cdot {\alpha} _{r_j}[j] \wedge (1 +{\alpha} _{r_j}[j])=0 \in \Lambda ^{2} k_{\infty} ^{\times}.   
 $$
 
\end{lemma}

\begin{theorem}\label{theorem-0-identity} Suppose that we are given a $\nu$-periodic sequence of mutations in a cluster pattern as in (\ref{periodic-sequence}). Let  $k$ be a field with ${\rm char}(k)=0,$   and  $\alpha_i \in k_{m} ,$ for $1\leq i \leq n,$ such that   for every $0 \leq j <P,$ the corresponding $\alpha _{r_j}[j]$ has the property that $-\alpha_{r_j}[j]  \in k_{m} ^{\flat}.$  Then  we have  
$$
\sum _{0\leq j <P} \theta_{r_j}\cdot \ell i _{m,w}(-\alpha_{r_j}[j])=0.
 $$

    
\end{theorem}

\begin{proof}  

Let $R$ be the polynomial ring over  $\mathbb{Q}$ generated by the indeterminates $x _{i,e},$ with $1 \leq i \leq n$ and $0\leq e <m,$ let $F$ be the field of fractions of $R$   and $x_i := \sum _{0 \leq e <m} x _{i,e}t^e \in F_m.$ Applying the mutations above appearing in the $\nu$-periodic sequence, we obtain $x_{i}[j] \in F_{m},$ for $1\leq i\leq n$ and $0\leq j<P.$  In order to ease the notation, we put $y_{i}[j]:=x _{i}[j](0)$ and $y_i:=y_{i}[0].$

Let us also put
$$
f(x_{i,e})_{{1\leq i \leq n} \atop {0\leq e <m}} :=   \sum _{0\leq j <P} \theta_{r_j}\cdot \ell i _{m,w}(-x_{r_j}[j]) \in F.
$$
Notice that $f$ is a rational function in the variables $x_{i,e}$ with coefficients in $\mathbb{Q}.$  From the definition of $\ell i_{m,w},$ we see that $f$  has poles only along some irreducible polynomials in $\mathbb{Q}[y_1,\cdots, y_n] \subseteq R.$ 
If $f$ has a pole along an irreducible polynomial $p(y_1,\cdots,y_n)$ in $\mathbb{Q}[y_1,\cdots, y_n]$ then 
there is a $j$  such that the valuation of $y_{r_j}[j]$ or $1+y_{r_j}[j]$ at $p(y_1,\cdots, y_n)$ is non-zero. Denote by $Y$ the algebraic subset of $\mathbb{A}^{n}_{\mathbb{Q}}$ defined by the product of those irreducible polynomials $p(y_{1},\cdots, y_{n})$  such that there is a $j$  with the property that the valuation of $y_{r_j}[j]$ or $1+y_{r_j}[j]$ at $p(y_1,\cdots, y_n)$ is non-zero. We then  have $Y \subseteq X \subseteq \mathbb{A}^{n} _{\mathbb{Q}},$ with $X$ as in Lemma \ref{lemma}. Note that for $\alpha _{i} \in k_{m},$ with $1 \leq i \leq n,$ the condition that $-\alpha _{r_j}[j] \in k_{m} ^{\flat}$ for every $0\leq j<P$ is equivalent to  $(\alpha_{1}(0), \cdots, \alpha _{n}(0) ) \in (\mathbb{A}^{n}_{\mathbb{Q}} \setminus Y)(k)$ and in this case, 
\begin{align}\label{rational-eqn}
  f(\alpha_{i,e})_{{1\leq i \leq n} \atop {0\leq e <m}} =   \sum _{0\leq j <P} \theta_{r_j}\cdot \ell i _{m,w}(-\alpha_{r_j}[j]) \in k,  
\end{align}
where we put $\alpha_{i}=\sum _{0\leq e<m}\alpha_{i,e}t^e.$



Let $\overline{k}$ denote an algebraic closure of $k.$  For $\beta _{i} \in \overline{k}_m$ such that $(\beta_1(0),\cdots, \beta_{n}(0)) \in (\mathbb{A} ^n \setminus X)(\overline{k}),$ we choose  $\tilde{\beta}_{i} \in \overline{k}_{\infty} $ which reduce to  $\beta _i$ modulo $(t^m).$ By  Lemma \ref{lemma}, we have
\begin{eqnarray}\label{theta-eqn}
  \sum_{0\leq j <P}\theta _{r_{j}}\cdot (1+\tilde{\beta} _{r_j}[j]) \wedge \tilde{\beta} _{r_j}[j]=0.     
\end{eqnarray}

 
For any  $\beta \in \overline{k}_m ^{\flat}$ and any $\tilde{\beta} \in \overline{k}_{\infty} ^{\flat},$ which  reduces to  $\beta $ modulo $(t^m),$ we have, by (\ref{dilog-formula}),
\begin{align}\label{lifting-indentity-0}
\ell i _{m,w}(-\beta)=\sum_{1 \leq i \leq w - m} i \cdot (\ell_{w - i} \land \ell_i)((1+\tilde{\beta}) \wedge \tilde{\beta}).    
\end{align}
This implies that  
$$
\sum _{0\leq j <P} \theta_{r_j}\ell i _{m,w}(-\beta_{r_j}[j])=\sum _{0\leq j <P} \sum _{1\leq i \leq w-m}\theta_{r_j} i \cdot (\ell_{w - i} \land \ell_i)((1+\tilde{\beta}_{r_j}[j]) \wedge \tilde{\beta}_{r_j}[j]).
$$
The right hand side can be rewritten as 
$$
 \sum _{1\leq i \leq w-m}i \cdot (\ell_{w - i} \land \ell_i)\Big ( \sum _{0\leq j <P}\theta_{r_j}\cdot ((1+\tilde{\beta}_{r_j}[j]) \wedge \tilde{\beta}_{r_j}[j])\Big).
$$
We have shown above that the sum in parentheses is equal to 0 in (\ref{theta-eqn}). By (\ref{rational-eqn}), this implies that $f(\beta_{i,e})_{{1\leq i \leq n} \atop {0\leq e <m}} =0,$ for $\beta_{i} \in \overline{k}_{m}$ such that $(\beta_1(0),\cdots, \beta_{n}(0)) \in (\mathbb{A} ^n \setminus X)(\overline{k}).$ Since $\overline{k}$ is algebraically closed and  $f$ is a rational function that does not have poles along $Y, $ this implies that $f(\beta_{i,e})_{{1\leq i \leq n} \atop {0\leq e <m}} =0,$ for $\beta_{i} \in \overline{k}_{m}$ with $(\beta_1(0),\cdots, \beta_{n}(0)) \in (\mathbb{A} ^n \setminus Y)(\overline{k}).$ Using (\ref{rational-eqn}) shows that 
$$
\sum _{0\leq j <P} \theta_{r_j}\cdot \ell i _{m,w}(-\beta_{r_j}[j])=0,
$$
if $(\beta_{1}(0),\cdots, \beta_{n}(0)) \in (\mathbb{A}^{n}\setminus Y)(\overline{k}).$ Since $k\subseteq \overline{k},$ we have the statement in the theorem. 
\end{proof}

\end{document}
